\begin{abstract}
Automatic screening of phonocardiogram (PCG) recordings could extend cardiac
auscultation to settings without expert cardiology, but on the PhysioNet/CinC
2016 corpus a high aggregate score is easy to obtain and hard to trust, because
the recording site is itself strongly predictive of the diagnosis. This project
compares three feature representations (MFCC, log-Mel spectrogram, perceptual
wavelet packet) and three classifiers (SVM, random forest, small CNN) under one
recording-level protocol, and pairs each performance claim with a diagnostic
that could refute it: unsupervised $k$-means structure, a per-site breakdown,
and a falsifiable test of whether attributions concentrate in systole. The best
model reaches a modified accuracy of $0.866$, comparable to published entries,
but $k$-means recovers the acquisition site far better than the diagnosis, and
per-site modified accuracy ranges from $0.99$ down to $0.48$. A direct test
confirms the mechanism: the recording site is recoverable from the same
features at $0.65$ balanced accuracy (versus $0.17$ chance), and diagnostic performance falls
from $0.84$ to $0.46$ --- essentially chance --- under leave-one-site-out
evaluation, though a weak signal survives once site is held fixed. Attribution is
only weakly systole-enriched, and no more so for correct than for incorrect
predictions. We conclude that the aggregate metric substantially overstates what the models
have learned.
\end{abstract}
